\scanpage{109}

\mspage{115}

\diagram{\begin{tikzcd}[ampersand replacement=\&,row sep=small]
 \mathbb D\text{-équiv. }2\text{-complète}\arrow[d,Rightarrow] \& {}\\
 \mathbb D\text{-}\unsure{\text{équi}|\text{équiv.}}\ 1\text{-complète}\arrow[r,Leftrightarrow] \& \mathbb D\text{-équiv. coh.}\arrow[d,Rightarrow]\\
 \& \mathbb D\text{-équiv. hom.}\arrow[d,Rightarrow]\\
 \& \mathbb D\text{-équivalence}
\end{tikzcd}}

Pour \deleted{\unsure{HOT}} $\mathbb D=\mathrm{HOT}$, toutes \unsure{les|ces} notions devraient être équivalentes, et équivalentes \deleted{en fait} à la condition encore plus faible en apparence, que $f:X\to Y$ soit une Hot-équivalence.

Mais je ne connais que HOT comme dérivateur où la $\mathbb D$-équivalence d'une flèche $f:X\to Y$ de $\mathrm{Diag}$ implique déjà la $\mathbb D$-équiv. coh. (complète), ou la $\mathbb D$-équivalence homologique (complète), voire même la $\mathbb D$-équivalence $1$-complète. (Je reviendrai là-dessus tantôt.)

\uline{Prop. 1} Soit $f:X\to Y$ de $\mathrm{Diag}$. Conditions équivalentes :

a) Pour tout $y\in\mathrm{Ob}\,Y$, l'inclusion $X_y\to X_{/y}$ est une $\mathbb D$-équiv., et pour toute flèche $y'\to y$ de $Y$, $X_{/y'}\to\unsure{X_{/y}|Y_{/y}}$ est une $\mathbb D$-équiv.

b) Pour tout \deleted{faisceau}\inserted{coeff.} \uline{constant} $F$ sur $Y$, $f_* f^* F$ est loc. constant, et se calcule ``fibre par fibre''.

(C'est la prop. 1 p. \unsure{76|75})

\uline{Prop. 1'} duale

% archive sheet number: 109
